\documentclass[11pt]{article}

% -------------------------------------------------------------------
% Standalone preregistration note. Self-contained: no dependency on the
% Second-Edition macros, so it preregisters on its own. It EXTRACTS and
% QUOTES the frozen register (rd:s:frozen) and Appendix G of QGT r26.9;
% it does not alter the frozen record.
% -------------------------------------------------------------------

\usepackage[a4paper,margin=1in]{geometry}
\usepackage{amsmath,amssymb}
\usepackage{booktabs}
\usepackage[hidelinks]{hyperref}
\usepackage{enumitem}

\newcommand{\vp}{\varphi}
\newcommand{\ainv}{\alpha^{-1}}

\title{\bfseries A Frozen Recoil--Branch Prediction\\ for the Fine-Structure Constant}
\author{Pasquale Camelia\\
\normalsize ORCID: 0009-0006-4779-4647}
\date{Frozen 4 July 2026 --- preregistration extract from QGT Second Edition r26.9}

\begin{document}
\maketitle

\begin{abstract}
\noindent This note isolates, for preregistration, a single quantitative
prediction of Quantum Gauge Theory (QGT): the values a photon-recoil
determination of the fine-structure constant must return. The prediction is
signed, dated ahead of the data that will test it, and falsifiable. It is
extracted verbatim from the frozen section of the Second Edition (r26.9,
\texttt{rd:s:frozen}) and from Register~B of that volume's Prediction Ledger
(Appendix~G); the underlying theory is not re-derived here, and none of it needs
to hold for this note to be tested. What is claimed, what is admitted as input,
what would confirm it, and what falls if it fails are all stated on this page. Its
register, in the language of the volume, is a \emph{conditional frozen forecast}:
prospective and dated, but conditional on one discrete choice made \emph{a
posteriori} on existing data, and declared as such. It is not offered as an
operatorially derived prediction, and it is not offered as proof of the theory.
\end{abstract}

\section{The frozen values}

Against the parameter-free fixed point $\ainv_{\rm fp}=137.0359991703$, the
frozen prediction for the two improved recoil determinations is

\begin{center}
\small
\renewcommand{\arraystretch}{1.3}\setlength{\tabcolsep}{4pt}
\begin{tabular}{lll}
\toprule
Determination & frozen prediction (4 July 2026) & current value\\
\midrule
$^{87}$Rb recoil, improved & $\ainv=137.035999215$ & $137.035999206(11)$\ \ [Morel 2020]\\
$^{133}$Cs recoil, improved & $\ainv=137.035999053$ & $137.035999046(27)$\ \ [Parker 2018]\\
any new recoil species (Sr, Yb,\ldots) & $w\in\{-0.7236,\,0,\,+0.2764\}$ & ---\\
\bottomrule
\end{tabular}
\end{center}

\noindent The recoil weight $w_A$ of a species is defined by the apparatus
contraction $\ainv_A=\ainv_{\rm fp}+w_A\,f_1$, with the single-channel rate
$f_1=\vp\cdot10^{-7}=1.618\times10^{-7}$. The blind content of the third row is
its \emph{discreteness}: a species not yet measured is forced onto three signed
weights, not a continuum. No atomic datum enters the construction of the lattice
itself.

\section{Construction and admitted inputs}

The values rest on three theorems of the volume meeting one hypothesis; the note
states them only to expose the inputs, not to prove them.

\begin{itemize}[leftmargin=1.4em,itemsep=2pt]
\item \textbf{Theorems (fixed structure, no atomic input).} The apparatus
contraction $\ainv_A=\ainv_{\rm fp}+w_A f_1$; the branch complementarity of the
Schur return, giving the visible weight $Z_\vp=\vp/\sqrt5\approx0.7236$ and the
hidden weight $1-Z_\vp=1/(\sqrt5\,\vp)\approx0.2764$; and the residue reading of
couplings. From these the branch displacements have exact closed forms in the
golden ring,
\[
\Delta_{\rm hidden}=+\frac{10^{-7}}{\sqrt5},\qquad
\Delta_{\rm visible}=-\frac{\vp^2\,10^{-7}}{\sqrt5},\qquad
\Delta_{\rm hidden}-\Delta_{\rm visible}=\vp\cdot10^{-7}=f_1 .
\]

\item \textbf{Hypothesis.} A channel-weighted readout couples to \emph{one}
branch of the return, so admissible weights are the discrete lattice
$w_A\in\{-Z_\vp,\,0,\,1-Z_\vp\}$, with self-readouts at $w=0$.

\item \textbf{The one a-posteriori choice, declared.} Which measured channel
reads which branch was assigned using the two existing recoil data: rubidium on
the hidden branch, caesium on the visible branch with opposite orientation.
Because that selection used the data, the $0.06\,\sigma$ agreement of the
measured splitting with $f_1$ is part of the selection, \emph{not} an independent
check on past data. This is why the register is conditional, and it is the sharp
line between this note and a derived prediction.
\end{itemize}

\noindent What is prospective, fixed with no further parameter once the branch
pair is chosen: the splitting equals $f_1$ exactly (the branch weights differ by
one), and the asymmetry is
$|w_{\rm Cs}|/|w_{\rm Rb}|=Z_\vp/(1-Z_\vp)=\vp^2$. These stand as predictions on
data yet to come.

\section{Current status: survival, not confirmation}

With the branch assignment above, both existing determinations fit: rubidium at
$0.82\,\sigma$ on the hidden branch, caesium at $0.27\,\sigma$ on the visible
branch. Among the golden assignments considered, each rival is dead beyond
$2.3\,\sigma$ on at least one face --- equipartition $\pm\tfrac12$ at
$4.1\,\sigma$, the pair $(\vp^{-2},-\vp^{-1})$ at $2.4\,\sigma$, the pair
$(0,-1)$ at $3.3\,\sigma$, all on the rubidium face. Two disciplines are kept:
the rival list is the list \emph{considered}, not a predefined exhaustive space,
so no global exclusion and no look-elsewhere correction is claimed; and, as
above, the splitting agreement is inherited from the a-posteriori selection. The
recoil pair therefore \emph{survives}. Survival is not confirmation; the blind
test is the next improved determination.

\section{The decisive test, with thresholds}

The prediction is signed, and the signs are its exposure. The true rubidium value
is predicted to sit \emph{above} the current central value, by
$+9\times10^{-9}$ on $\ainv$; the true caesium value likewise slightly above its
current central value, by $+7\times10^{-9}$. The test becomes decisive at three
standard deviations once a recoil determination reaches an uncertainty near
$3\times10^{-9}$ on $\ainv$ --- a relative precision of $2\times10^{-11}$, one
further factor of four beyond the 2020 rubidium measurement, on the trajectory
that programme is already on. No new apparatus class is required; the test is a
precision improvement on an existing measurement.

\section{Negative controls and the competing null}

The competing null explanation is retained explicitly and is not treated as
excluded: a conventional unrecognised systematic in one interferometer would
reproduce a Rb--Cs offset without any of the present structure, and the
metrological community holds exactly this possibility open. The reading here is
\emph{not} a spacetime variation of $\alpha$: the invariant does not drift and
only the channel contraction depends on the apparatus, so existing atomic-clock
limits on a temporal drift of $\alpha$ are consistent with it and do not test it.
The negative control for the branch structure is convergence: if improved
determinations move together toward one value, the structure that this note
predicts is absent. A $^{229}$Th nuclear clock is deliberately \emph{excluded}
from the recoil lattice: it reads a sensitivity $\delta\nu/\nu=K_\alpha\,
\delta\alpha/\alpha+\cdots$, not an absolute $h/M$ recoil comparison, and placing
it on the lattice would compare inhomogeneous observables. It enters elsewhere as
a distinct, doubly conditional target ($K_{\rm Th}$ and $w_{\rm Th}$ to be
derived first), named as a target and \emph{not} predicted here.

\section{Falsifiers, and what falls}

The falsifiers are enumerated so that no retreat is available afterwards.

\begin{enumerate}[leftmargin=1.6em,itemsep=2pt]
\item \textbf{Convergence.} If improved recoil determinations converge to a single
common value well inside $10^{-7}$, the channel-weighted class collapses and the
tensorial reading of the coupling fails.
\item \textbf{Off-lattice weight.} If any new recoil weight lands off the lattice
$\{-0.7236,\,0,\,+0.2764\}$ at three standard deviations, the branch hypothesis
fails even if the tensorial reading survives.
\item \textbf{Sign failure.} If the improved rubidium value drifts \emph{away}
from $137.035999215$ at three standard deviations, the assignment fails in sign
--- the same death.
\end{enumerate}

\noindent What falls if it fails is bounded and stated: a failure of (2) or (3)
kills the branch assignment and this note, and leaves the rest of the volume
standing; a failure of (1) is heavier --- it removes the channel-weighted reading
of $\alpha$ on which this note and the readout spectrum both rest. Neither outcome
requires, or refutes, the remainder of QGT. That containment is the point of
preregistering the recoil claim on its own.

\section*{Promotion path}

One open step would turn this conditional forecast into a derived prediction: the
channel-to-branch map from the measurement operator --- a demonstration, from the
contraction an $h/M$-against-$R_\infty$ comparison actually performs, of
\emph{why} one recoil protocol reads the hidden branch and the other the visible
branch with opposite orientation, sign included. That map is the single missing
piece; the frozen numbers above do not wait for it.

\section*{Freeze and reproducibility}

The values were constructed on 4~July~2026, ahead of any atomic-recoil or
nuclear-clock determination of $\alpha$ published after that date, and are frozen.
The machine-checkable record recomputes every frozen number from the volume's own
formulas (never hard-coded) and recomputes the current-data distances and rival
exclusions from the primary 2018/2020 inputs.

\begin{center}
\small
\begin{tabular}{ll}
\toprule
Frozen record & \texttt{tools/qgt\_branch\_prediction\_frozen.py}\\
Record SHA-256 (unchanged since r24.1) & \texttt{fb68dbbb869f5ade\ldots56ce36c1}\\
r26.x source-tree SHA-256 & \texttt{14a7a68be8aac900\ldots b8220e08}\\
\bottomrule
\end{tabular}
\end{center}

\noindent After deposit in an immutable archive with a DOI, the frozen register
may change only to record a \textsc{pass}, \textsc{fail}, or \textsc{inconclusive}
verdict of the data --- never its numerical content.

\bigskip
\footnotesize\sloppy
\noindent\textit{Sources (QGT Second Edition r26.9):} \S\,``The frozen
prediction: the branch lattice of the recoil weights'' (\texttt{rd:s:frozen},
ch.~08); ``The Prediction Ledger,'' Register~A and Register~B (Appendix~G);
\texttt{FREEZE\_MANIFEST.txt}. \textit{Data:} Morel et al., \textit{Nature}
\textbf{588}, 61 (2020) [Rb]; Parker et al., \textit{Science} \textbf{360}, 191
(2018) [Cs].

\end{document}
